\chapter{Gram-Schmidt and QR}

There is a recurring theme in this book: \emph{perpendicular is easier}. Whenever
two directions are at right angles, the bookkeeping simplifies. Lengths add by
Pythagoras instead of by the law of cosines. Projections decouple, one coordinate
at a time, instead of tangling together in a linear system. Coordinates can be read
off by a single dot product instead of solved for.

Suppose you are given a set of vectors that span a subspace---the columns of a matrix,
for example---but they point in oblique directions.
Could you replace this basis with an \textbf{orthonormal basis}---a set of vectors
that are mutually orthogonal ($90^\circ$) and each of unit length---spanning
the same subspace and constructed directly from the original vectors?

The procedure to do this is \textbf{Gram--Schmidt}, which means
``project away what you already have, keep what is new.'' When you record the
arithmetic Gram--Schmidt performs, that record \emph{is} a factorization of your
original matrix: $\mat{A} = \mat{Q}\mat{R}$, with $\mat{Q}$ having orthonormal
columns and $\mat{R}$ \textbf{upper triangular} (every entry below the main
diagonal is zero). This is the \textbf{QR factorization}, a standard tool in
numerical linear algebra and stable least-squares solvers.

\section{The wish: straighten a crooked basis}

Consider two vectors in the plane, $\vec{a}_1$ and $\vec{a}_2$, that are linearly independent
but not orthogonal (Figure~\ref{ch20:crooked}). They span $\R^2$, so they form a valid
basis. However, finding the coordinates of a vector $\vec{b}$ with respect to this basis
requires solving a $2\times 2$ linear system, because the directions are coupled: a
vector along $\vec{a}_1$ has a nonzero component along $\vec{a}_2$.

\begin{figure}[H]
\centering
\begin{tikzpicture}[scale=1.3,>=Stealth]
  \draw[->,gray!50] (-0.4,0)--(3.2,0) node[right,font=\small]{$x$};
  \draw[->,gray!50] (0,-0.4)--(0,2.6) node[above,font=\small]{$y$};
  \draw[->,very thick,vecblue] (0,0)--(2.6,0.6) node[right]{$\vec{a}_1$};
  \draw[->,very thick,vecred]  (0,0)--(1.5,2.0) node[above]{$\vec{a}_2$};
  % Angle arc between a1 and a2
  \draw (13:0.65) arc (13:53:0.65);
  \node at (33:0.85){\small $\theta$};
  % Shadow of a2 on line of a1
  \draw[dashed,spacegray] (1.5,2.0)--(1.95,0.45);
  \draw[->,thick,structurepurple] (0,0)--(1.95,0.45)
    node[midway,below=2pt,font=\footnotesize]{\small shadow of $\vec{a}_2$ on $\vec{a}_1$};
\end{tikzpicture}
\caption{A non-orthogonal basis. The vectors $\vec{a}_1,\vec{a}_2$ span the plane but meet
at an angle $\theta \neq 90^\circ$. Projecting $\vec{a}_2$ onto $\vec{a}_1$
yields a nonzero component: the directions are coupled, so finding coordinates requires
solving a system. An orthonormal basis avoids this coupling.}
\label{ch20:crooked}
\end{figure}

\begin{recoverybox}
\textbf{Pause and think.} If two unit vectors $\vec{q}_1$ and $\vec{q}_2$ are
perpendicular, what is the length of their sum $\vec{q}_1+\vec{q}_2$? By Pythagoras,
$\norm{\vec{q}_1+\vec{q}_2}^2 = \norm{\vec{q}_1}^2 + \norm{\vec{q}_2}^2 = 1+1=2$, so
the length is exactly $\sqrt{2}$. With crooked vectors, cross-terms appear and spoil
this simplicity. Perpendicularity removes cross-terms from every calculation.
\end{recoverybox}

\begin{intuitionbox}
In an orthonormal basis, every vector has unit length and distinct vectors are mutually
orthogonal. Coordinates are then computed directly by dot products: the component of $\vec{b}$
along $\vec{q}_i$ is simply $\inner{\vec{q}_i}{\vec{b}}$, without solving a linear
system. Gram--Schmidt is the algorithm that converts any basis into an orthonormal
basis spanning the same space.
\end{intuitionbox}

Recall from the chapter on projections that the projection of a vector $\vec{a}$ onto
the line through a vector $\vec{q}$ is
\[
  \proj_{\vec{q}}(\vec{a}) \;=\; \frac{\inner{\vec{q}}{\vec{a}}}{\inner{\vec{q}}{\vec{q}}}\,\vec{q}.
\]
If $\vec{q}$ is already a unit vector, the denominator is $1$ and this simplifies to
$\inner{\vec{q}}{\vec{a}}\,\vec{q}$. That single formula is the only tool we need.

\section{Gram-Schmidt for two vectors}

Start with independent vectors $\vec{a}_1$ and $\vec{a}_2$. We want an orthonormal
pair $\vec{q}_1,\vec{q}_2$ spanning the same space.

\textbf{Step 1: normalize the first vector.} There are no preceding directions to satisfy
orthogonality against, so normalize $\vec{a}_1$ to unit length:
\[
  \vec{q}_1 \;=\; \frac{\vec{a}_1}{\norm{\vec{a}_1}}.
\]

\textbf{Step 2: remove the part of $\vec{a}_2$ that lies along $\vec{q}_1$.}
The vector $\vec{a}_2$ has a shadow on the line of $\vec{q}_1$ -- the component that
$\vec{q}_1$ already covers. Subtract it off. What remains is exactly the part of
$\vec{a}_2$ that points in a genuinely \emph{new} direction, perpendicular to
$\vec{q}_1$:
\[
  \vec{v}_2 \;=\; \vec{a}_2 - \proj_{\vec{q}_1}(\vec{a}_2)
            \;=\; \vec{a}_2 - \inner{\vec{q}_1}{\vec{a}_2}\,\vec{q}_1.
\]
Then normalize:
\[
  \vec{q}_2 \;=\; \frac{\vec{v}_2}{\norm{\vec{v}_2}}.
\]

\begin{figure}[H]
\centering
\begin{tikzpicture}[scale=1.5,>=Stealth]
  \draw[->,gray!40] (-0.3,0)--(3.0,0) node[right,font=\small]{$x$};
  \draw[->,gray!40] (0,-0.3)--(0,2.3) node[above,font=\small]{$y$};
  % a1 along x-axis
  \draw[->,line width=1.4pt,vecblue!50] (0,0)--(2.5,0) node[below right,font=\small]{$\vec{a}_1$};
  % unit vector q1
  \draw[->,very thick,vecblue] (0,0)--(1.0,0) node[above=2pt,font=\small]{$\vec{q}_1$};
  % a2 leaning
  \draw[->,very thick,vecred] (0,0)--(1.5,1.7) node[above]{$\vec{a}_2$};
  % projection (shadow)
  \draw[->,thick,structurepurple] (0,0)--(1.5,0);
  \node[structurepurple,below=2pt,font=\footnotesize] at (0.75,0){$\proj_{\vec{q}_1}(\vec{a}_2)$};
  % perpendicular leftover v2
  \draw[->,very thick,targetgreen] (1.5,0)--(1.5,1.7)
    node[midway,right=3pt,font=\small]{$\vec{v}_2 \perp \vec{q}_1$};
  % right-angle marker
  \draw (1.35,0)--(1.35,0.15)--(1.5,0.15);
  % unit vector q2 on y-axis
  \draw[->,very thick,targetgreen!80!black] (0,0)--(0,1.0)
    node[left=2pt,font=\small]{$\vec{q}_2$};
  \fill (0,0) circle (1.2pt) node[below left,font=\tiny]{$\vec{0}$};
\end{tikzpicture}
\caption{Gram--Schmidt on two vectors. The first direction $\vec{a}_1$ normalizes to $\vec{q}_1$.
The second vector $\vec{a}_2$ casts a shadow $\proj_{\vec{q}_1}(\vec{a}_2) = \inner{\vec{q}_1}{\vec{a}_2}\vec{q}_1$.
Subtracting that shadow leaves $\vec{v}_2$, which is perpendicular to $\vec{q}_1$. Normalizing
$\vec{v}_2$ gives $\vec{q}_2$, yielding the clean orthonormal pair $(\vec{q}_1,\vec{q}_2)$.}
\label{ch20:gs2}
\end{figure}

Why is $\vec{v}_2$ guaranteed perpendicular to $\vec{q}_1$? Take the dot product and
watch it cancel. Using $\inner{\vec{q}_1}{\vec{q}_1}=1$,
\[
  \inner{\vec{q}_1}{\vec{v}_2}
   = \inner{\vec{q}_1}{\vec{a}_2} - \inner{\vec{q}_1}{\vec{a}_2}\,\inner{\vec{q}_1}{\vec{q}_1}
   = \inner{\vec{q}_1}{\vec{a}_2} - \inner{\vec{q}_1}{\vec{a}_2}
   = 0.
\]
This cancellation happens because we subtracted off \emph{exactly} the amount that
overlapped $\vec{q}_1$, so nothing of $\vec{q}_1$ is left in $\vec{v}_2$.

\begin{recoverybox}
\textbf{Pause and think.} What if $\vec{a}_1$ and $\vec{a}_2$ were already perpendicular
at the start? Then $\inner{\vec{q}_1}{\vec{a}_2}=0$, so the shadow is $\vec{0}$. The
subtraction leaves $\vec{v}_2 = \vec{a}_2$ unchanged, and normalizing produces
$\vec{q}_2 = \vec{a}_2/\norm{\vec{a}_2}$. Gram--Schmidt leaves already-perpendicular
vectors completely undisturbed, adjusting only their lengths.
\end{recoverybox}

\begin{labbox}
Let $\vec{a}_1=\colvec{3\\0}$ and $\vec{a}_2=\colvec{2\\2}$.
\begin{itemize}
  \item $\norm{\vec{a}_1}=3$, so $\vec{q}_1=\colvec{1\\0}$.
  \item $\inner{\vec{q}_1}{\vec{a}_2}=2$, so the shadow is $2\,\vec{q}_1=\colvec{2\\0}$.
  \item $\vec{v}_2=\colvec{2\\2}-\colvec{2\\0}=\colvec{0\\2}$, and
        $\vec{q}_2=\colvec{0\\1}$.
\end{itemize}
Check: $\inner{\vec{q}_1}{\vec{q}_2}=0$, both unit length. We recovered the standard
axes -- which makes sense, since $\vec{a}_1$ was already on the $x$-axis. Try it again
with $\vec{a}_1=\colvec{1\\1}$ and see the frame come out tilted by $45^\circ$.
\end{labbox}

\section{Gram-Schmidt in general}

The pattern now repeats. When we reach $\vec{a}_k$, we already have an orthonormal set
$\vec{q}_1,\dots,\vec{q}_{k-1}$. We subtract off \emph{all} the shadows at once -- the
component of $\vec{a}_k$ along each $\vec{q}_i$ we have built so far -- and whatever
survives is perpendicular to every one of them.

\begin{formalbox}
\textbf{The Gram--Schmidt process.} Given independent vectors
$\vec{a}_1,\dots,\vec{a}_n$, define for $k=1,2,\dots,n$:
\[
  \vec{v}_k \;=\; \vec{a}_k \;-\; \sum_{i=1}^{k-1} \inner{\vec{q}_i}{\vec{a}_k}\,\vec{q}_i,
  \qquad
  \vec{q}_k \;=\; \frac{\vec{v}_k}{\norm{\vec{v}_k}}.
\]
Then $\vec{q}_1,\dots,\vec{q}_n$ is an orthonormal set with
$\spn\{\vec{q}_1,\dots,\vec{q}_k\}=\spn\{\vec{a}_1,\dots,\vec{a}_k\}$ for every $k$.
\end{formalbox}

The summation subtracts the projection of $\vec{a}_k$ onto each previously constructed
direction $\vec{q}_i$. The remainder $\vec{v}_k$ is the component of $\vec{a}_k$ that is
orthogonal to $\spn\{\vec{q}_1,\dots,\vec{q}_{k-1}\}$.

\begin{focusbox}
\textbf{Project, subtract, normalize.} At step $k$, project $\vec{a}_k$ onto the
subspace spanned by the earlier directions, subtract that projection to isolate the
orthogonal component $\vec{v}_k$, and normalize $\vec{v}_k$ to unit length.
\end{focusbox}

Let us verify both claims -- orthogonality and equal spans -- directly from the definition.

\begin{proposition}[Orthogonality]
\label{ch20:orth-prop}
For each $k$, the vector $\vec{v}_k$ defined above is orthogonal to every
$\vec{q}_j$ with $j<k$.
\end{proposition}

\begin{proof}
Fix $j<k$ and dot $\vec{q}_j$ against $\vec{v}_k$:
\[
  \inner{\vec{q}_j}{\vec{v}_k}
   = \inner{\vec{q}_j}{\vec{a}_k}
     - \sum_{i=1}^{k-1}\inner{\vec{q}_i}{\vec{a}_k}\,\inner{\vec{q}_j}{\vec{q}_i}.
\]
By induction the earlier vectors $\vec{q}_1,\dots,\vec{q}_{k-1}$ are already
orthonormal, so $\inner{\vec{q}_j}{\vec{q}_i}=0$ except when $i=j$, where it equals
$1$. The sum therefore collapses to its single $i=j$ term:
\[
  \inner{\vec{q}_j}{\vec{v}_k}
   = \inner{\vec{q}_j}{\vec{a}_k} - \inner{\vec{q}_j}{\vec{a}_k}\cdot 1 = 0. \qedhere
\]
\end{proof}

\begin{proposition}[Equal spans]
\label{ch20:span-prop}
For each $k$, $\spn\{\vec{q}_1,\dots,\vec{q}_k\}=\spn\{\vec{a}_1,\dots,\vec{a}_k\}$.
\end{proposition}

\begin{proof}
Each $\vec{q}_k$ is built only from $\vec{a}_k$ and the earlier
$\vec{q}_1,\dots,\vec{q}_{k-1}$ -- which in turn came from $\vec{a}_1,\dots,\vec{a}_{k-1}$
-- so every $\vec{q}_k$ is a combination of $\vec{a}_1,\dots,\vec{a}_k$. Conversely,
rearranging the definition gives
\[
  \vec{a}_k \;=\; \norm{\vec{v}_k}\,\vec{q}_k \;+\; \sum_{i=1}^{k-1}\inner{\vec{q}_i}{\vec{a}_k}\,\vec{q}_i,
\]
so every $\vec{a}_k$ is a combination of $\vec{q}_1,\dots,\vec{q}_k$. Each spanning set
sits inside the other; the spans are equal. Here $\norm{\vec{v}_k}\neq 0$ because the
$\vec{a}_k$ are independent, so no projection ever lands exactly on $\vec{a}_k$.
\end{proof}

\begin{trapbox}
\textbf{Normalize as you go, or keep the denominator.} The clean projection formula
$\inner{\vec{q}_i}{\vec{a}_k}\,\vec{q}_i$ assumes each finished vector $\vec{q}_i$ has
unit length. If you prefer to delay normalizing and work with the unnormalized vectors
$\vec{v}_i$ instead, you must keep the denominator
$\frac{\inner{\vec{v}_i}{\vec{a}_k}}{\inner{\vec{v}_i}{\vec{v}_i}}\,\vec{v}_i$ in every
projection term. Both formulations produce identical orthogonal directions; forgetting
to divide by $\norm{\vec{v}_i}^2$ when skipping normalization is the classic student bug.
\end{trapbox}

\begin{walkbox}
You have built $\vec{q}_1,\dots,\vec{q}_{k-1}$ and you are about to process
$\vec{a}_k$. Before computing anything: what does the leftover $\vec{v}_k$ \emph{mean}
geometrically? (Hint: stack the $\vec{q}_i$ into a subspace and think about the
perpendicular from $\vec{a}_k$ to that subspace -- it is the same projection picture
from the projections chapter, just onto a span instead of a single line.)
\end{walkbox}

Figure~\ref{ch20:gsplane} makes that perpendicular concrete. The shaded plane is the
span of the directions already finished; the sum of shadows
$\sum_i \inner{\vec{q}_i}{\vec{a}_k}\,\vec{q}_i$ is the foot of the perpendicular from
$\vec{a}_k$, and $\vec{v}_k$ is the perpendicular itself -- the part of $\vec{a}_k$ that
sticks out of the old subspace.

\begin{figure}[H]
\centering
\begin{tikzpicture}[scale=1.0,>=Stealth,line join=round]
  \tikzset{x={(1cm,0cm)},y={(0.5cm,0.32cm)},z={(0cm,1cm)}}
  % the already-built subspace as a shaded parallelogram in the x-y "floor"
  \fill[sourceblue!12] (-0.4,-0.4,0)--(3.0,-0.4,0)--(3.0,2.2,0)--(-0.4,2.2,0)--cycle;
  \node[sourceblue] at (2.6,2.0,0){\small $\spn\{\vec{q}_1,\dots,\vec{q}_{k-1}\}$};
  % two finished directions lying in the floor
  \draw[->,very thick,vecblue] (0,0,0)--(1.6,0,0) node[anchor=north]{$\vec{q}_1$};
  \draw[->,very thick,vecblue] (0,0,0)--(0,1.6,0) node[anchor=west]{$\vec{q}_{k-1}$};
  % a_k sticking up out of the plane
  \draw[->,very thick,vecred] (0,0,0)--(1.3,1.1,1.5) node[anchor=south west]{$\vec{a}_k$};
  % its foot (the sum of shadows) in the plane
  \draw[->,thick,structurepurple] (0,0,0)--(1.3,1.1,0);
  \node[structurepurple] at (1.5,0.55,0){\small $\sum_i\inner{\vec{q}_i}{\vec{a}_k}\vec{q}_i$};
  % the perpendicular leftover v_k
  \draw[->,very thick,targetgreen] (1.3,1.1,0)--(1.3,1.1,1.5);
  \node[targetgreen] at (1.75,1.1,0.8){$\vec{v}_k$};
  % right-angle hint
  \draw (1.15,0.95,0)--(1.15,0.95,0.18)--(1.3,1.1,0.18);
\end{tikzpicture}
\caption{Gram--Schmidt in general is projection onto a \emph{plane} (or higher
subspace), not just a line. The foot of the perpendicular is the sum of the shadows on
the finished directions; the leftover $\vec{v}_k$ stands straight up out of
$\spn\{\vec{q}_1,\dots,\vec{q}_{k-1}\}$, then gets normalized to $\vec{q}_k$.}
\label{ch20:gsplane}
\end{figure}

\section{A fully worked three-vector example}

We trace the procedure on three vectors in $\R^3$. Take
\[
  \vec{a}_1=\colvec{1\\1\\0},\qquad
  \vec{a}_2=\colvec{1\\0\\1},\qquad
  \vec{a}_3=\colvec{0\\1\\1}.
\]

\textbf{First direction.} $\norm{\vec{a}_1}=\sqrt{1^2+1^2+0^2}=\sqrt2$, so
\[
  \vec{q}_1 = \frac{1}{\sqrt2}\colvec{1\\1\\0}.
\]

\textbf{Second direction.} First the overlap with $\vec{q}_1$:
\[
  \inner{\vec{q}_1}{\vec{a}_2}
   = \frac{1}{\sqrt2}(1\cdot 1 + 1\cdot 0 + 0\cdot 1) = \frac{1}{\sqrt2}.
\]
Subtract that shadow from $\vec{a}_2$:
\[
  \vec{v}_2
   = \colvec{1\\0\\1} - \frac{1}{\sqrt2}\cdot\frac{1}{\sqrt2}\colvec{1\\1\\0}
   = \colvec{1\\0\\1} - \frac12\colvec{1\\1\\0}
   = \colvec{\tfrac12\\[2pt]-\tfrac12\\[2pt]1}.
\]
Its length is $\norm{\vec{v}_2}=\sqrt{\tfrac14+\tfrac14+1}=\sqrt{\tfrac32}=\dfrac{\sqrt6}{2}$, so
\[
  \vec{q}_2 = \frac{2}{\sqrt6}\colvec{\tfrac12\\[2pt]-\tfrac12\\[2pt]1}
            = \frac{1}{\sqrt6}\colvec{1\\-1\\2}.
\]
Quick check: $\inner{\vec{q}_1}{\vec{q}_2}=\frac{1}{\sqrt2}\frac{1}{\sqrt6}(1-1+0)=0$.
Perpendicular, as promised.

\textbf{Third direction.} Now subtract \emph{both} shadows from $\vec{a}_3$. Compute
the two overlaps:
\[
  \inner{\vec{q}_1}{\vec{a}_3}=\frac{1}{\sqrt2}(0+1+0)=\frac{1}{\sqrt2},
  \qquad
  \inner{\vec{q}_2}{\vec{a}_3}=\frac{1}{\sqrt6}(0-1+2)=\frac{1}{\sqrt6}.
\]
Then
\[
  \vec{v}_3
   = \colvec{0\\1\\1}
     - \frac{1}{\sqrt2}\cdot\frac{1}{\sqrt2}\colvec{1\\1\\0}
     - \frac{1}{\sqrt6}\cdot\frac{1}{\sqrt6}\colvec{1\\-1\\2}.
\]
The coefficients are $\tfrac12$ and $\tfrac16$, so
\[
  \vec{v}_3
   = \colvec{0\\1\\1}
     - \frac12\colvec{1\\1\\0}
     - \frac16\colvec{1\\-1\\2}
   = \colvec{0-\tfrac12-\tfrac16 \\[2pt] 1-\tfrac12+\tfrac16 \\[2pt] 1-0-\tfrac13}
   = \colvec{-\tfrac23\\[2pt]\tfrac23\\[2pt]\tfrac23}.
\]
Its length is $\norm{\vec{v}_3}=\sqrt{3\cdot\tfrac49}=\sqrt{\tfrac43}=\dfrac{2}{\sqrt3}$, so
\[
  \vec{q}_3 = \frac{\sqrt3}{2}\colvec{-\tfrac23\\[2pt]\tfrac23\\[2pt]\tfrac23}
            = \frac{1}{\sqrt3}\colvec{-1\\1\\1}.
\]
Final sanity check on the third vector against the first two:
\[
  \inner{\vec{q}_1}{\vec{q}_3}=\tfrac{1}{\sqrt6}(-1+1+0)=0,
  \qquad
  \inner{\vec{q}_2}{\vec{q}_3}=\tfrac{1}{\sqrt{18}}(-1-1+2)=0.
\]
Three mutually orthogonal unit vectors, constructed from the original basis.

\begin{figure}[H]
\centering
\begin{tikzpicture}[scale=0.95,>=Stealth]
  % ===== PANEL 1: STEP 1 =====
  \begin{scope}[shift={(-5.0,0)}]
    \draw[fill=gray!2, draw=gray!25, rounded corners=3pt] (-2.2,-2.3) rectangle (2.2,2.85);
    \node[draw=sourceblue!70, fill=sourceblue!12, rounded corners=2pt, font=\footnotesize\scshape, text=sourceblue!90!black]
      at (0,2.5) {Step 1: Normalize $\vec{a}_1$};

    \draw[->,spacegray] (-1.8,0)--(1.8,0) node[right,font=\tiny]{$x$};
    \draw[->,spacegray] (0,-1.8)--(0,1.8) node[above,font=\tiny]{$y$};

    % a1 and q1
    \draw[->,line width=1.6pt,vecblue!50] (0,0)--(1.6,1.6) node[above,font=\scriptsize]{$\vec{a}_1$};
    \draw[->,very thick,vecblue] (0,0)--(1.0,1.0) node[below right,font=\scriptsize]{$\vec{q}_1$};
    \fill (0,0) circle (1.6pt) node[below left,font=\tiny]{$\vec{0}$};

    \node[font=\tiny, spacegray, align=center] at (0,-1.6) {
      $\vec{q}_1 = \dfrac{\vec{a}_1}{\norm{\vec{a}_1}}$\\[2pt]
      Establish baseline unit axis
    };
  \end{scope}

  % ===== PANEL 2: STEP 2 =====
  \begin{scope}[shift={(0,0)}]
    \draw[fill=gray!2, draw=gray!25, rounded corners=3pt] (-2.2,-2.3) rectangle (2.2,2.85);
    \node[draw=targetgreen!70, fill=targetgreen!12, rounded corners=2pt, font=\footnotesize\scshape, text=targetgreen!90!black]
      at (0,2.5) {Step 2: Strip Shadow};

    \draw[->,spacegray] (-1.8,0)--(1.8,0) node[right,font=\tiny]{$x$};
    \draw[->,spacegray] (0,-1.8)--(0,1.8) node[above,font=\tiny]{$y$};

    % q1 along diagonal
    \draw[->,thick,vecblue] (0,0)--(1.0,1.0) node[above left,font=\tiny]{$\vec{q}_1$};
    % a2
    \draw[->,line width=1.4pt,vecred!70] (0,0)--(1.5,0) node[below,font=\scriptsize]{$\vec{a}_2$};
    % shadow of a2 on q1
    \draw[->,thick,structurepurple] (0,0)--(0.75,0.75) node[midway,above left=2pt,font=\tiny,fill=white,inner sep=1pt]{$\proj$};
    % drop dashed line
    \draw[dashed,gray] (1.5,0)--(0.75,0.75);
    % v2 perpendicular
    \draw[->,very thick,targetgreen!85!black] (0.75,0.75)--(1.5,0);
    \node[targetgreen!85!black,font=\tiny,right] at (1.2,0.45){$\vec{v}_2 \perp \vec{q}_1$};
    % normalized q2
    \draw[->,very thick,targetgreen] (0,0)--(0.7,-0.7) node[below left,font=\scriptsize]{$\vec{q}_2$};
    \fill (0,0) circle (1.6pt);

    \node[font=\tiny, spacegray, align=center] at (0,-1.6) {
      $\vec{v}_2 = \vec{a}_2 - \inner{\vec{q}_1}{\vec{a}_2}\vec{q}_1$\\[2pt]
      $\vec{q}_2 = \vec{v}_2 / \norm{\vec{v}_2} \perp \vec{q}_1$
    };
  \end{scope}

  % ===== PANEL 3: STEP 3 =====
  \begin{scope}[shift={(5.0,0)}]
    \draw[fill=gray!2, draw=gray!25, rounded corners=3pt] (-2.2,-2.3) rectangle (2.2,2.85);
    \node[draw=quotientorange!70, fill=quotientorange!12, rounded corners=2pt, font=\footnotesize\scshape, text=quotientorange!90!black]
      at (0,2.5) {Step 3: Plane Remnant};

    % 3D floor plane spanned by q1, q2
    \fill[sourceblue!10] (-1.6,-0.9) -- (1.5,-0.9) -- (0.8,0.4) -- (-1.9,0.4) -- cycle;
    \draw[sourceblue!50,thin] (-1.6,-0.9) -- (1.5,-0.9) -- (0.8,0.4) -- (-1.9,0.4) -- cycle;
    \node[sourceblue!70!black,font=\tiny] at (-1.1,0.2) {$\spn\{\vec{q}_1,\vec{q}_2\}$};

    \draw[->,thick,vecblue] (0,-0.4)--(0.8,-0.4) node[below,font=\tiny]{$\vec{q}_1$};
    \draw[->,thick,targetgreen] (0,-0.4)--(-0.6,-0.1) node[left,font=\tiny]{$\vec{q}_2$};

    % a3 rising up
    \draw[->,line width=1.4pt,vecred] (0,-0.4)--(0.7,1.4) node[above,font=\scriptsize]{$\vec{a}_3$};
    % projection in floor
    \fill[structurepurple] (0.7,-0.2) circle (1.4pt);
    \draw[dashed,gray] (0.7,1.4)--(0.7,-0.2);
    \draw[dashed,structurepurple] (0,-0.4)--(0.7,-0.2);
    % v3 vertical
    \draw[->,line width=1.5pt,quotientorange!95!black] (0.7,-0.2)--(0.7,1.4)
      node[midway,right=1pt,font=\tiny]{$\vec{v}_3$};
    % normalized q3 at origin
    \draw[->,very thick,quotientorange!90!black] (0,-0.4)--(0,0.8)
      node[left=1pt,font=\scriptsize]{$\vec{q}_3$};
    \fill (0,-0.4) circle (1.6pt);

    \node[font=\tiny, spacegray, align=center] at (0,-1.6) {
      $\vec{v}_3 = \vec{a}_3 - \sum_{i=1}^2\inner{\vec{q}_i}{\vec{a}_3}\vec{q}_i$\\[2pt]
      $\vec{q}_3 \perp \{\vec{q}_1, \vec{q}_2\}$
    };
  \end{scope}
\end{tikzpicture}
\caption{The three sequential steps of Gram--Schmidt orthogonalization.
\emph{Step 1:} normalize the first vector $\vec{a}_1$ to set the initial unit direction $\vec{q}_1$.
\emph{Step 2:} project $\vec{a}_2$ onto $\vec{q}_1$ and strip off that shadow; the leftover $\vec{v}_2$
is orthogonal to $\vec{q}_1$ and normalizes to $\vec{q}_2$.
\emph{Step 3:} project $\vec{a}_3$ onto the plane $\spn\{\vec{q}_1,\vec{q}_2\}$, strip off both shadows,
leaving the vertical remnant $\vec{v}_3$ that stands perpendicular to the entire plane; normalize to $\vec{q}_3$.}
\label{ch20:gs3}
\end{figure}

\section{The QR factorization}

The steps of the Gram--Schmidt process directly yield a matrix factorization.
Rearranging the formula from Proposition~\ref{ch20:span-prop}:
\[
  \vec{a}_k = \inner{\vec{q}_1}{\vec{a}_k}\,\vec{q}_1
            + \dots
            + \inner{\vec{q}_{k-1}}{\vec{a}_k}\,\vec{q}_{k-1}
            + \norm{\vec{v}_k}\,\vec{q}_k.
\]
Every original vector $\vec{a}_k$ is written as a combination of
$\vec{q}_1,\dots,\vec{q}_k$ -- strictly \emph{only} those, never the later
$\vec{q}$'s, because $\vec{q}_{k+1},\dots$ did not exist yet when we processed
$\vec{a}_k$. That ``only the earlier $\vec{q}$'s'' fact is exactly what makes a matrix
\emph{upper triangular}.

Stack the $\vec{a}_k$ as the columns of a matrix $\mat{A}$ and the $\vec{q}_k$ as the
columns of $\mat{Q}$. The coefficients above fill a matrix $\mat{R}$, and the relation
becomes a single clean equation:
\[
  \boxed{\;\mat{A} = \mat{Q}\mat{R}\;}
\]
where $\mat{Q}$ has orthonormal columns and $\mat{R}$ is upper triangular with the
entries
\[
  R_{ik} = \inner{\vec{q}_i}{\vec{a}_k}\ \ (i<k),
  \qquad
  R_{kk} = \norm{\vec{v}_k},
  \qquad
  R_{ik}=0\ \ (i>k).
\]

\begin{intuitionbox}
The factorization $\mat{A}=\mat{Q}\mat{R}$ separates geometry from coordinates. The matrix
$\mat{Q}$ contains the orthonormal basis vectors, while the upper-triangular matrix
$\mat{R}$ contains the coordinates that express each original column in this basis.
Reading column $k$ of $\mat{A}=\mat{Q}\mat{R}$ expresses $\vec{a}_k$ as a linear
combination of $\vec{q}_1,\dots,\vec{q}_k$.
\end{intuitionbox}

The triangular structure of $\mat{R}$ reflects the order in which Gram--Schmidt proceeds:
each $\vec{a}_k$ depends only on $\vec{q}_1,\dots,\vec{q}_k$. The diagonal entries
$R_{kk}=\norm{\vec{v}_k}$ are strictly positive whenever the columns $\vec{a}_k$ are
linearly independent.

\begin{figure}[H]
\centering
\begin{tikzpicture}[scale=1.15,>=Stealth]
  \draw[->,spacegray] (-0.3,0)--(4.2,0) node[right,font=\small]{$\vec{q}_1$};
  \draw[->,spacegray] (0,-0.3)--(0,2.3) node[above,font=\small]{$\vec{q}_2$};
  \fill[sourceblue!12] (0,0)--(2.8,0)--(3.8,2.0)--(1.0,2.0)--cycle;
  \draw[sourceblue!60,thick] (0,0)--(2.8,0)--(3.8,2.0)--(1.0,2.0)--cycle;
  \draw[->,line width=1.6pt,vecblue] (0,0)--(2.8,0) node[above right,font=\small]{$\vec{a}_1$};
  \draw[->,line width=1.6pt,vecred] (0,0)--(1.0,2.0) node[above left,font=\small]{$\vec{a}_2$};
  \draw[dashed,spacegray] (1.0,0)--(1.0,2.0);
  \draw (0.85,0)--(0.85,0.15)--(1.0,0.15);
  \draw[->,very thick,black] (0,0)--(1.0,0) node[midway,above=1pt,font=\footnotesize]{$\vec{q}_1$};
  \draw[->,very thick,black] (0,0)--(0,1.0) node[midway,left=1pt,font=\footnotesize]{$\vec{q}_2$};
  \draw[|<->|,vecblue] (0,-0.45)--(2.8,-0.45)
    node[midway,below,font=\footnotesize]{$R_{11} = \norm{\vec{a}_1}$ (base)};
  \draw[|<->|,structurepurple] (0,2.5)--(1.0,2.5)
    node[midway,above,font=\footnotesize]{$R_{12} = \inner{\vec{q}_1}{\vec{a}_2}$ (slant)};
  \draw[dashed,spacegray] (2.8,0)--(3.9,0);
  \draw[dashed,spacegray] (3.8,2.0)--(3.9,2.0);
  \draw[|<->|,targetgreen!80!black] (3.9,0)--(3.9,2.0)
    node[midway,right,font=\footnotesize]{$R_{22} = \norm{\vec{v}_2}$ (altitude)};
  \node[sourceblue!80!black,align=center,font=\small] at (1.9,1.0)
    {Parallelogram Area\\$= R_{11}R_{22} = \det\mat{R}$};
  \fill (0,0) circle (1.2pt) node[below left,font=\tiny]{$\vec{0}$};
\end{tikzpicture}
\caption{The geometric meaning of $\mat{A}=\mat{Q}\mat{R}$. In the orthonormal coordinates
given by $(\vec{q}_1,\vec{q}_2)$, the first column $\vec{a}_1$ lies along $\vec{q}_1$ with
base length $R_{11}=\norm{\vec{a}_1}$. The second column $\vec{a}_2$ slants along $\vec{q}_1$
by $R_{12}=\inner{\vec{q}_1}{\vec{a}_2}$ and reaches perpendicular altitude $R_{22}=\norm{\vec{v}_2}$
along $\vec{q}_2$. The area of the spanned parallelogram is $\text{base}\times\text{height} = R_{11}R_{22}$,
which is precisely the determinant $\det\mat{R}$.}
\label{ch20:qr-geom}
\end{figure}

\begin{recoverybox}
\textbf{Pause and think.} Why must every diagonal entry $R_{kk}=\norm{\vec{v}_k}$ be
strictly positive whenever the columns of $\mat{A}$ are linearly independent? If
$\norm{\vec{v}_k}=0$, then $\vec{v}_k=\vec{0}$, which would mean $\vec{a}_k = \sum_{i<k}
\inner{\vec{q}_i}{\vec{a}_k}\vec{q}_i$. In words: $\vec{a}_k$ would be a linear
combination of the earlier vectors, contradicting independence! Linear independence
guarantees $R_{kk}>0$, ensuring $\mat{R}$ is invertible.
\end{recoverybox}

\subsection{QR for our three-vector example}

We already have everything. The dot products we computed are precisely the entries of
$\mat{R}$. With
\[
  \mat{Q}=
  \begin{bmatrix}
    \tfrac{1}{\sqrt2} & \tfrac{1}{\sqrt6} & -\tfrac{1}{\sqrt3}\\[4pt]
    \tfrac{1}{\sqrt2} & -\tfrac{1}{\sqrt6} & \tfrac{1}{\sqrt3}\\[4pt]
    0 & \tfrac{2}{\sqrt6} & \tfrac{1}{\sqrt3}
  \end{bmatrix},
  \qquad
  \mat{R}=
  \begin{bmatrix}
    \norm{\vec{a}_1} & \inner{\vec{q}_1}{\vec{a}_2} & \inner{\vec{q}_1}{\vec{a}_3}\\[4pt]
    0 & \norm{\vec{v}_2} & \inner{\vec{q}_2}{\vec{a}_3}\\[4pt]
    0 & 0 & \norm{\vec{v}_3}
  \end{bmatrix}
  =
  \begin{bmatrix}
    \sqrt2 & \tfrac{1}{\sqrt2} & \tfrac{1}{\sqrt2}\\[4pt]
    0 & \tfrac{\sqrt6}{2} & \tfrac{1}{\sqrt6}\\[4pt]
    0 & 0 & \tfrac{2}{\sqrt3}
  \end{bmatrix}.
\]
We check the columns of $\mat{Q}\mat{R}$ directly. Column
$1$ of $\mat{R}$ is $(\sqrt2,0,0)\T$, so
\[
  \mat{Q}\,\colvec{\sqrt2\\0\\0}
  = \sqrt2\,\vec{q}_1
  = \sqrt2\cdot\frac{1}{\sqrt2}\colvec{1\\1\\0}
  = \colvec{1\\1\\0}
  = \vec{a}_1. \checkmark
\]
And column $2$: $\mat{R}$'s second column is
$(\tfrac{1}{\sqrt2},\tfrac{\sqrt6}{2},0)\T$, giving
\[
  \frac{1}{\sqrt2}\vec{q}_1 + \frac{\sqrt6}{2}\vec{q}_2
  = \frac12\colvec{1\\1\\0} + \frac{\sqrt6}{2}\cdot\frac{1}{\sqrt6}\colvec{1\\-1\\2}
  = \frac12\colvec{1\\1\\0} + \frac12\colvec{1\\-1\\2}
  = \colvec{1\\0\\1}
  = \vec{a}_2. \checkmark
\]

\begin{labbox}
Factor $\mat{A}=\begin{bmatrix}1&2\\0&1\\1&0\end{bmatrix}$ by hand. The columns are
$\vec{a}_1=(1,0,1)\T$ and $\vec{a}_2=(2,1,0)\T$.
\begin{itemize}
  \item $\norm{\vec{a}_1}=\sqrt2$, $\vec{q}_1=\tfrac{1}{\sqrt2}(1,0,1)\T$.
  \item $\inner{\vec{q}_1}{\vec{a}_2}=\tfrac{1}{\sqrt2}(2+0+0)=\sqrt2$.
  \item $\vec{v}_2=(2,1,0)\T-\sqrt2\cdot\tfrac{1}{\sqrt2}(1,0,1)\T=(1,1,-1)\T$,
        $\norm{\vec{v}_2}=\sqrt3$, $\vec{q}_2=\tfrac{1}{\sqrt3}(1,1,-1)\T$.
\end{itemize}
Then $\mat{R}=\begin{bmatrix}\sqrt2 & \sqrt2\\ 0 & \sqrt3\end{bmatrix}$. Multiply
$\mat{Q}\mat{R}$ back out and confirm you recover $\mat{A}$.
\end{labbox}

\subsection{Properties of orthogonal matrices}

A matrix whose columns are orthonormal satisfies the single most useful identity in
this corner of the subject:
\[
  \mat{Q}\T\mat{Q} = \mat{I}.
\]
Entry $(i,j)$ of $\mat{Q}\T\mat{Q}$ is $\inner{\vec{q}_i}{\vec{q}_j}$, which is $1$ when
$i=j$ and $0$ otherwise. When $\mat{Q}$ is square it is called \textbf{orthogonal},
and its inverse is simply its transpose: $\mat{Q}\inv=\mat{Q}\T$.
Orthogonal matrices preserve lengths and angles: $\norm{\mat{Q}\vec{x}}=\norm{\vec{x}}$
and $\inner{\mat{Q}\vec{x}}{\mat{Q}\vec{y}}=\inner{\vec{x}}{\vec{y}}$. Geometrically,
multiplication by a square orthogonal matrix represents a rigid motion (a rotation or
reflection). Computationally, preserving lengths prevents the amplification of numerical errors.

\begin{strangbox}
Strang introduces orthonormal bases and Gram--Schmidt alongside projections and
least squares, expressing $\mat{A}=\mat{Q}\mat{R}$ in this form. The foundational
geometric idea is that of Figure~\ref{ch20:gs2}: project, subtract, normalize. The next
sections show how $\mat{Q}\T\mat{Q}=\mat{I}$ simplifies the normal equations into a
triangular solve, and why the modified algorithm is used in numerical software.
\end{strangbox}

\section{Why orthonormal columns make least squares easy}

Recall the least-squares problem from the previous chapter: $\mat{A}\vec{x}=\vec{b}$
usually has no solution because $\vec{b}$ is not in the column space, so we settle for
the $\hat{\vec{x}}$ that minimizes $\norm{\mat{A}\vec{x}-\vec{b}}$. The classic route
is the \textbf{normal equations}
\[
  \mat{A}\T\mat{A}\,\hat{\vec{x}} = \mat{A}\T\vec{b}.
\]
This system is mathematically exact, but explicitly forming $\mat{A}\T\mat{A}$ can
amplify roundoff errors by squaring the \textbf{condition number} of $\mat{A}$.
The condition number $\kappa(\mat{A})$ measures sensitivity to perturbations: a
relative error of size $\eps$ in the data can produce a relative error of order
$\kappa(\mat{A})\eps$ in the computed solution. Forming $\mat{A}\T\mat{A}$ squares this
quantity: $\kappa(\mat{A}\T\mat{A}) = \kappa(\mat{A})^2$. For example, if $\kappa(\mat{A}) = 10^4$
(costing about 4 decimal digits of precision), $\mat{A}\T\mat{A}$ has condition
number $10^8$, losing 8 digits of precision before solving begins. The QR factorization
avoids forming $\mat{A}\T\mat{A}$ directly.

Substituting $\mat{A}=\mat{Q}\mat{R}$ into the normal equations simplifies the system.
First,
\[
  \mat{A}\T\mat{A} = (\mat{Q}\mat{R})\T(\mat{Q}\mat{R})
                   = \mat{R}\T\mat{Q}\T\mat{Q}\mat{R}
                   = \mat{R}\T\mat{R},
\]
using $\mat{Q}\T\mat{Q}=\mat{I}$. Likewise
$\mat{A}\T\vec{b}=\mat{R}\T\mat{Q}\T\vec{b}$. The normal equations become
$\mat{R}\T\mat{R}\,\hat{\vec{x}}=\mat{R}\T\mat{Q}\T\vec{b}$. Cancel the common
$\mat{R}\T$ (it is invertible, since its diagonal entries $\norm{\vec{v}_k}$ are
nonzero) to land at the strikingly simple
\[
  \boxed{\;\mat{R}\,\hat{\vec{x}} = \mat{Q}\T\vec{b}.\;}
\]

\begin{focusbox}
Least squares via QR proceeds in two steps: compute the vector $\mat{Q}\T\vec{b}$,
then solve the upper-triangular system $\mat{R}\hat{\vec{x}}=\mat{Q}\T\vec{b}$ by
back-substitution. This avoids forming $\mat{A}\T\mat{A}$ and preserves the condition
number $\kappa(\mat{A})$.
\end{focusbox}

There is a clean geometric reading too. The projection of $\vec{b}$ onto the column
space of $\mat{A}$ -- the best approximation $\hat{\vec{b}}$ -- is simply
$\hat{\vec{b}}=\mat{Q}\mat{Q}\T\vec{b}$, where $\mat{Q}\mat{Q}\T$ is the projection
matrix onto $\col(\mat{A})$. The orthonormal columns let you read off the projection as
a sum of independent one-dimensional projections,
$\hat{\vec{b}}=\sum_i \inner{\vec{q}_i}{\vec{b}}\,\vec{q}_i$, where each coordinate is
computed independently without coupling across directions (Figure~\ref{ch20:lsproj}).

\begin{figure}[H]
\centering
\begin{tikzpicture}[scale=1.0,>=Stealth,line join=round]
  \tikzset{x={(1cm,0cm)},y={(0.5cm,0.32cm)},z={(0cm,1cm)}}
  % column space = plane spanned by q1,q2
  \fill[sourceblue!12] (-0.4,-0.4,0)--(3.0,-0.4,0)--(3.0,2.2,0)--(-0.4,2.2,0)--cycle;
  \node[sourceblue] at (2.55,2.0,0){\small $\col(\mat{A})$};
  \draw[->,very thick,vecblue] (0,0,0)--(1.5,0,0) node[anchor=north]{$\vec{q}_1$};
  \draw[->,very thick,vecblue] (0,0,0)--(0,1.5,0) node[anchor=west]{$\vec{q}_2$};
  % data vector b out of the plane
  \draw[->,very thick,vecred] (0,0,0)--(1.2,1.0,1.4) node[anchor=south west]{$\vec{b}$};
  % its projection onto the column space
  \draw[->,thick,structurepurple] (0,0,0)--(1.2,1.0,0);
  \node[structurepurple] at (1.5,0.5,0){\small $\hat{\vec{b}}=\mat{Q}\mat{Q}\T\vec{b}$};
  % the residual (error) perpendicular to the plane
  \draw[->,very thick,targetgreen] (1.2,1.0,0)--(1.2,1.0,1.4);
  \node[targetgreen] at (1.7,1.0,0.75){\small $\vec{b}-\hat{\vec{b}}$};
  \draw (1.05,0.85,0)--(1.05,0.85,0.18)--(1.2,1.0,0.18);
\end{tikzpicture}
\caption{Least squares is one projection. The best approximation $\hat{\vec{b}}$ is the
shadow of $\vec{b}$ on the column space, $\hat{\vec{b}}=\mat{Q}\mat{Q}\T\vec{b}$, and the
residual $\vec{b}-\hat{\vec{b}}$ stands perpendicular to that space. Orthonormal columns
allow computing each coordinate of $\hat{\vec{b}}$ independently.}
\label{ch20:lsproj}
\end{figure}

\begin{recoverybox}
\textbf{Pause and think.} When $\mat{Q}$ is rectangular ($m \times n$ with $m>n$),
$\mat{Q}\T\mat{Q} = \mat{I}_n$ is the $n\times n$ identity, but $\mat{Q}\mat{Q}\T$ is
an $m\times m$ matrix that is \emph{not} the identity. What does $\mat{Q}\mat{Q}\T$ do?
It is the orthogonal projector onto $\col(\mat{A})$: it leaves any vector already in the
column space alone, and zeroes out any vector perpendicular to it.
\end{recoverybox}

\begin{example}[Least squares by QR]
\label{ch20:ls-example}
Fit a line through the three points $(0,1),(1,2),(2,2)$ -- the same kind of fit from
the regression chapter. We want $c+dx\approx y$, so
\[
  \mat{A}=\begin{bmatrix}1&0\\1&1\\1&2\end{bmatrix},\qquad
  \vec{b}=\colvec{1\\2\\2}.
\]
Gram--Schmidt on the columns gives
$\vec{q}_1=\tfrac{1}{\sqrt3}(1,1,1)\T$ and, after subtracting,
$\vec{v}_2=(0,1,2)\T-1\cdot(1,1,1)\T=(-1,0,1)\T$, so
$\vec{q}_2=\tfrac{1}{\sqrt2}(-1,0,1)\T$. Here
$\inner{\vec{q}_1}{\vec{a}_2}=\tfrac{1}{\sqrt3}(0+1+2)=\sqrt3$ and
$\norm{\vec{v}_2}=\sqrt2$, so
\[
  \mat{R}=\begin{bmatrix}\sqrt3 & \sqrt3\\ 0 & \sqrt2\end{bmatrix},
  \qquad
  \mat{Q}\T\vec{b}=
  \colvec{\tfrac{1}{\sqrt3}(1+2+2)\\[4pt]\tfrac{1}{\sqrt2}(-1+0+2)}
  =\colvec{\tfrac{5}{\sqrt3}\\[4pt]\tfrac{1}{\sqrt2}}.
\]
Back-substitute $\mat{R}\hat{\vec{x}}=\mat{Q}\T\vec{b}$. The bottom row gives
$\sqrt2\,d=\tfrac{1}{\sqrt2}$, so $d=\tfrac12$. The top row gives
$\sqrt3\,c+\sqrt3\,d=\tfrac{5}{\sqrt3}$, i.e. $c+d=\tfrac53$, so
$c=\tfrac53-\tfrac12=\tfrac76$. The best line is $y=\tfrac76+\tfrac12 x$, matching what
the normal equations would give -- but reached by a triangular solve.
\end{example}

\begin{applicationbox}
\textbf{QR inside every regression and every deep-net training loop.} When you call
\texttt{numpy.linalg.lstsq}, fit an ordinary least-squares model in a statistics
package, or solve the linear sub-problem inside a Gauss--Newton or
Levenberg--Marquardt optimizer, the library does \emph{not} form $\mat{A}\T\mat{A}$.
It computes a QR factorization (or the closely related SVD) and back-substitutes,
exactly because squaring the condition number can destroy accuracy on real,
near-collinear data. In large-scale machine learning, QR also underlies stable
orthogonalization steps: ``QR-based whitening,'' re-orthogonalizing the basis vectors
in Krylov methods like \textsc{gmres} and Lanczos (the iterative solvers behind huge
sparse systems), and orthogonal weight initialization, where a network's weight matrix
is set to the $\mat{Q}$ of a random matrix so that signals neither explode nor vanish
as they pass through layers. Whenever preserving orthogonality and vector lengths is
essential, QR factorization is standard.
\end{applicationbox}

\begin{connectionbox}
\textbf{Gram--Schmidt depends only on the inner product, not on $\R^n$.} The
process used \emph{only} dot products, lengths, and projections. Swap the dot product
for any other inner product and the identical recipe orthogonalizes those objects. The
famous case: on the space of polynomials with the inner product
$\inner{f}{g}=\int_{-1}^{1} f(x)g(x)\,dx$, applying Gram--Schmidt to the plain powers
$1,x,x^2,x^3,\dots$ produces the \textbf{Legendre polynomials}, the orthogonal basis at
the heart of Gaussian quadrature and spherical harmonics. Other weightings give
Chebyshev, Hermite, and Laguerre families. In this way, Gram--Schmidt provides a
unified procedure for constructing classical orthogonal polynomial systems in analysis.
\end{connectionbox}

\section{Classical versus modified Gram-Schmidt}

The version we derived is \textbf{classical} Gram--Schmidt: to build $\vec{v}_k$,
project the \emph{original} $\vec{a}_k$ onto each finished $\vec{q}_i$ and subtract all
the shadows from $\vec{a}_k$ at once. In exact arithmetic, this method is exact. In floating-point arithmetic, however,
computing each projection coefficient $\inner{\vec{q}_i}{\vec{a}_k}$ from the original
vector $\vec{a}_k$ allows roundoff errors to accumulate; when columns are nearly
dependent, the computed vectors can lose orthogonality.

\textbf{Modified Gram--Schmidt} addresses this instability by subtracting each projection
from the running remainder rather than from the original vector:

\begin{formalbox}
\textbf{Modified Gram--Schmidt.} Set $\vec{v}\leftarrow\vec{a}_k$. Then for
$i=1,2,\dots,k-1$ in order, update
\[
  \vec{v}\;\leftarrow\;\vec{v}-\inner{\vec{q}_i}{\vec{v}}\,\vec{q}_i,
\]
and finally $\vec{q}_k=\vec{v}/\norm{\vec{v}}$.
\end{formalbox}

\begin{figure}[H]
\centering
\begin{tikzpicture}[scale=0.95,>=Stealth]
  \begin{scope}[shift={(-3.8,0)}]
    \draw[fill=gray!2, draw=gray!25, rounded corners=3pt] (-3.2,-2.2) rectangle (3.2,2.0);
    \node[draw=sourceblue!70, fill=sourceblue!12, rounded corners=2pt, font=\footnotesize\bfseries, text=sourceblue!90!black]
      at (0,1.65) {Classical Gram--Schmidt (CGS)};
    \node[draw, thick, rounded corners, inner sep=4pt] (ak) at (-2.0, 0.2) {$\vec{a}_k$};
    \node[inner sep=2pt, font=\small] (p1) at (0.2, 0.8) {$\inner{\vec{q}_1}{\vec{a}_k}\vec{q}_1$};
    \node[inner sep=2pt, font=\small] (p2) at (0.2, -0.4) {$\inner{\vec{q}_2}{\vec{a}_k}\vec{q}_2$};
    \node[draw, thick, targetgreen!80!black, rounded corners, inner sep=4pt] (vk) at (2.2, 0.2) {$\vec{v}_k$};
    \draw[->, thick, spacegray] (ak) -- node[above, font=\tiny]{$\inner{\vec{q}_1}{\vec{a}_k}$} (p1);
    \draw[->, thick, spacegray] (ak) -- node[below, font=\tiny]{$\inner{\vec{q}_2}{\vec{a}_k}$} (p2);
    \draw[->, very thick, targetgreen] (p1) -- node[above right, font=\tiny]{subtract} (vk);
    \draw[->, very thick, targetgreen] (p2) -- node[below right, font=\tiny]{subtract} (vk);
    \node[font=\scriptsize, spacegray, align=center] at (0,-1.5) {
      All shadows projected from raw $\vec{a}_k$;\\
      early rounding errors do not get corrected.
    };
  \end{scope}

  \begin{scope}[shift={(3.8,0)}]
    \draw[fill=gray!2, draw=gray!25, rounded corners=3pt] (-3.2,-2.2) rectangle (3.2,2.0);
    \node[draw=targetgreen!70, fill=targetgreen!12, rounded corners=2pt, font=\footnotesize\bfseries, text=targetgreen!90!black]
      at (0,1.65) {Modified Gram--Schmidt (MGS)};
    \node[draw, thick, rounded corners, inner sep=4pt] (m0) at (-2.2, 0.2) {$\vec{a}_k$};
    \node[draw, thick, rounded corners, inner sep=4pt] (m1) at (0.0, 0.2) {$\vec{v}^{(1)}$};
    \node[draw, thick, targetgreen!80!black, rounded corners, inner sep=4pt] (m2) at (2.2, 0.2) {$\vec{v}_k$};
    \draw[->, very thick, structurepurple] (m0) -- node[above, font=\tiny]{$-\proj_{\vec{q}_1}$} (m1);
    \draw[->, very thick, targetgreen] (m1) -- node[above, font=\tiny]{$-\proj_{\vec{q}_2}$} (m2);
    \node[font=\scriptsize, spacegray, align=center] at (0,-1.5) {
      Next shadow dots against updated remainder;\\
      orthogonalization errors do not accumulate.
    };
  \end{scope}
\end{tikzpicture}
\caption{Classical versus Modified Gram--Schmidt. \emph{Classical:} all projections are
computed simultaneously from the original vector $\vec{a}_k$, so early errors are baked
into every term. \emph{Modified:} each shadow is subtracted sequentially, measuring overlap against the
updated remainder $\vec{v}^{(1)}$ to prevent error accumulation in floating-point arithmetic.}
\label{ch20:cgs-vs-mgs}
\end{figure}

In exact arithmetic this gives \emph{exactly} the same answer as the classical
version: subtracting components one at a time, each from the partially-cleaned vector,
removes the same total component. The difference appears in floating-point arithmetic.
After subtracting the $\vec{q}_1$-component, the remainder $\vec{v}$ is orthogonal to
$\vec{q}_1$ to within machine precision. Computing $\inner{\vec{q}_2}{\vec{v}}$
evaluates overlap against this updated remainder, preventing errors from previous
projections from compounding into subsequent steps.

\begin{trapbox}
The classical and modified algorithms differ by \emph{which vector you dot against}:
classical always uses the original $\vec{a}_k$; modified uses the running remainder
$\vec{v}$. In exact arithmetic, both algorithms yield identical vectors. In finite-precision
arithmetic, however, modified Gram--Schmidt maintains $\mat{Q}\T\mat{Q}$ much closer
to $\mat{I}$, making it the standard formulation in numerical code.
\end{trapbox}

\begin{remark}
Even modified Gram--Schmidt is not the gold standard for the worst-conditioned
problems. Production QR routines (\textsc{lapack}'s \texttt{dgeqrf}) instead build
$\mat{Q}$ from a sequence of \textbf{Householder reflections}---orthogonal transformations that reflect a vector
across a hyperplane to introduce zeros below the diagonal of an entire column at once.
Householder QR is backward stable and is the standard implementation in libraries such
as \texttt{numpy} and \textsc{matlab}. Gram--Schmidt remains conceptually fundamental
because it makes the underlying geometry explicit: repeated projection onto orthogonal
subspaces. We meet Householder reflections again in the numerical linear algebra chapters.
\end{remark}

\begin{historybox}
The algorithm we call Gram--Schmidt has a rich history driven by applied problems
across two centuries.

The earliest formulation came from Pierre-Simon Laplace in 1816. Working on celestial
mechanics -- calculating the gravitational orbits of Jupiter and Saturn so that minor
measurement errors would not compound over centuries -- Laplace wrote down an orthogonal
projection recipe to separate orbital parameters. In 1883, the Danish actuary and
mathematician J\o rgen Pedersen Gram independently derived the method to fit polynomial
curves to actuarial tables (life expectancy and pension records) by least squares; with
raw powers $1, x, x^2, \dots$, the normal equations were nearly singular and impossible
to solve reliably by hand. Twenty-four years later, in 1907, Erhard Schmidt reinvented
the procedure in G\"ottingen to construct orthonormal function bases for Fredholm
integral equations, where the unknown function appears inside an integral.

All three mathematicians worked with quill, paper, and printed logarithm tables. In
exact arithmetic, Gram--Schmidt is exact. When implemented on digital computers with
floating-point arithmetic, however, classical Gram--Schmidt suffers from numerical
instability: rounding errors during subtractions can compound rapidly, causing computed
columns to lose orthogonality.

The resolution came in two waves. First, numerical analysts introduced Modified
Gram--Schmidt, which projects against the running remainder rather than the raw original,
dramatically curbing error growth. Then in 1958, Alston Householder introduced elementary
orthogonal reflections -- reflecting entire columns onto coordinate axes at once using
full orthogonal matrices, which preserves orthogonality to the full precision of machine
hardware.

The naming reflects historical rediscovery: Gram's 1883 paper arose in actuarial mathematics,
while Schmidt's 1907 work applied to infinite-dimensional function spaces. Laplace's 1816
formulation was overlooked for more than a century before being recognized. Together, these
ideas led to the QR algorithm for eigenvalues (developed by John Francis and Vera
Kublanovskaya in 1961), modern least-squares solvers, and orthogonal polynomials throughout
quantum mechanics.
\end{historybox}

\begin{aestheticbox}
A single operation---subtracting orthogonal projections---yields an orthonormal basis
(Gram--Schmidt), a matrix factorization ($\mat{A}=\mat{Q}\mat{R}$), a stable method for
least squares ($\mat{R}\hat{\vec{x}}=\mat{Q}\T\vec{b}$), and families of orthogonal
polynomials under different inner products. The geometric principle of decomposing a
vector into orthogonal components provides both theoretical clarity and computational
stability throughout linear algebra.
\end{aestheticbox}

\section{Cost and intuition recap}

The procedure is summarized by three operations: \textbf{project, subtract, normalize}.
Key concepts:

\begin{itemize}
  \item \textbf{The wish.} A crooked basis tangles every coordinate computation into a
        linear solve. An orthonormal basis untangles it: each coordinate is one dot
        product, $\inner{\vec{q}_i}{\vec{b}}$.
  \item \textbf{The construction.} Gram--Schmidt builds the orthonormal frame one
        direction at a time. For each $\vec{a}_k$, subtract its shadows on the finished
        $\vec{q}_i$ and normalize what is left:
        $\vec{v}_k=\vec{a}_k-\sum_{i<k}\inner{\vec{q}_i}{\vec{a}_k}\vec{q}_i$,
        $\vec{q}_k=\vec{v}_k/\norm{\vec{v}_k}$.
  \item \textbf{The bookkeeping is a factorization.} Recording those dot products and
        lengths gives $\mat{A}=\mat{Q}\mat{R}$: $\mat{Q}$ holds clean directions,
        $\mat{R}$ (upper triangular, positive diagonal) holds the recipe.
  \item \textbf{The payoff.} $\mat{Q}\T\mat{Q}=\mat{I}$ turns the normal equations into
        the triangular solve $\mat{R}\hat{\vec{x}}=\mat{Q}\T\vec{b}$, dodging the
        condition-number squaring of $\mat{A}\T\mat{A}$.
\end{itemize}

\textbf{Cost.} For an $m\times n$ matrix ($m\ge n$), Gram--Schmidt (classical or
modified) costs about $2mn^2$ floating-point operations -- the same leading order as
forming and factoring $\mat{A}\T\mat{A}$, but with far better accuracy because the
errors never get squared. On a chalkboard the bottleneck is instead the square roots and
fractions in the normalizations; choosing friendly starting vectors keeps the arithmetic
clean. \textbf{Stability.} Modified Gram--Schmidt keeps $\mat{Q}\T\mat{Q}$ much closer to
$\mat{I}$ than classical; Householder QR (the production default) is better still. The
geometry, though, is identical across all three: subtract the projection.

\begin{connectionbox}
Geometrically, the $\mat{A} = \mat{Q}\mat{R}$ factorization also measures the volume of the parallelotope spanned by the columns of $\mat{A}$: the orthogonal matrix $\mat{Q}$ acts as a rigid transformation preserving volumes and angles, while the upper-triangular matrix $\mat{R}$ records lengths and perpendicular altitudes. The volume of the parallelotope is given by the product of these altitudes:
\[
  \text{Volume} \;=\; R_{11} R_{22} \cdots R_{nn} \;=\; \det\mat{R}.
\]
In Chapter~21 (``Determinants: Volume and Sign''), we examine how square matrices scale multi-dimensional volumes and how the sign of the determinant indicates spatial orientation. In Part~VI, we study eigenvalues and eigenvectors to understand how transformations behave along invariant directions. Later, in Chapter~29, orthogonal bases and spectral theory combine in the Singular Value Decomposition (SVD), a foundational tool of applied linear algebra.
\end{connectionbox}

\section{Exercises}

\begin{exercise}[Two-vector warmup]
Apply Gram--Schmidt to $\vec{a}_1=\colvec{2\\1}$ and $\vec{a}_2=\colvec{3\\4}$. Find
$\vec{q}_1,\vec{q}_2$ and verify they are orthonormal. Then write down $\mat{R}$ and
confirm $\mat{A}=\mat{Q}\mat{R}$.
\end{exercise}

\begin{exercise}[A clean three-vector run]
Orthonormalize
$\vec{a}_1=\colvec{1\\1\\1}$, $\vec{a}_2=\colvec{1\\1\\0}$, $\vec{a}_3=\colvec{1\\0\\0}$.
(They are independent, so the process will not break down.) Report
$\vec{q}_1,\vec{q}_2,\vec{q}_3$ and the matrix $\mat{R}$.
\end{exercise}

\begin{exercise}[Read off $\mat{R}$ without recomputing]
Suppose you have already produced $\vec{q}_1,\vec{q}_2$ from $\vec{a}_1,\vec{a}_2$ and
found $\norm{\vec{a}_1}=5$, $\inner{\vec{q}_1}{\vec{a}_2}=3$, $\norm{\vec{v}_2}=4$.
Write the full $2\times 2$ matrix $\mat{R}$. What is $\det\mat{R}$, and why does it
equal the area of the parallelogram spanned by $\vec{a}_1,\vec{a}_2$?
\end{exercise}

\begin{exercise}[Least squares by QR]
Fit a line to the points $(1,1),(2,2),(3,2),(4,5)$ using the QR route: build
$\mat{A},\vec{b}$, run Gram--Schmidt to get $\mat{Q},\mat{R}$, then solve
$\mat{R}\hat{\vec{x}}=\mat{Q}\T\vec{b}$ by back-substitution. Compare your slope and
intercept against solving the normal equations directly.
\end{exercise}

\begin{exercise}[Order matters for $\mat{Q}$, not for the span]
Run Gram--Schmidt on $\vec{a}_1=\colvec{1\\0\\0},\vec{a}_2=\colvec{1\\1\\0}$ and then
again on the \emph{reversed} list $\vec{a}_2,\vec{a}_1$. You get different orthonormal
bases. Explain geometrically why the two $\vec{q}_1$'s point differently but the spans
of $\{\vec{q}_1,\vec{q}_2\}$ are identical.
\end{exercise}

\begin{exercise}[Why $\mat{Q}\T\mat{Q}=\mat{I}$ but $\mat{Q}\mat{Q}\T\neq\mat{I}$]
Take the $\mat{Q}$ from Example~\ref{ch20:ls-example} (a $3\times 2$ matrix with
orthonormal columns). Compute both $\mat{Q}\T\mat{Q}$ and $\mat{Q}\mat{Q}\T$. One is
the $2\times 2$ identity; the other is a $3\times 3$ projection matrix. Identify which
subspace $\mat{Q}\mat{Q}\T$ projects onto, and explain why it cannot be the full
identity.
\end{exercise}

\begin{exercise}[Classical vs.\ modified, by hand]
Let $\eps$ be tiny and take
$\vec{a}_1=\colvec{1\\\eps\\0},\vec{a}_2=\colvec{1\\0\\\eps},\vec{a}_3=\colvec{1\\0\\0}$
(nearly dependent). Carry out both classical and modified Gram--Schmidt symbolically,
keeping terms to first order in $\eps$. Show that in exact arithmetic both give the same
$\vec{q}_3$, and describe in words where, with rounding, the classical version would let
a $\vec{q}_1$-error leak into $\vec{q}_3$ while the modified version would not.
\end{exercise}

\begin{exercise}[Orthogonal matrices preserve length -- proof]
Let $\mat{Q}$ be square with $\mat{Q}\T\mat{Q}=\mat{I}$. Prove that
$\norm{\mat{Q}\vec{x}}=\norm{\vec{x}}$ for every $\vec{x}$, and that
$\inner{\mat{Q}\vec{x}}{\mat{Q}\vec{y}}=\inner{\vec{x}}{\vec{y}}$ for all $\vec{x},
\vec{y}$. (Start from $\norm{\mat{Q}\vec{x}}^2=(\mat{Q}\vec{x})\T(\mat{Q}\vec{x})$.)
Conclude that multiplication by $\mat{Q}$ is a rotation or reflection, and connect this
to why QR-based computation does not amplify roundoff.
\end{exercise}

\begin{exercise}[Legendre by hand]
Using the inner product $\inner{f}{g}=\int_{-1}^{1}f(x)g(x)\,dx$, apply Gram--Schmidt
to $1,x,x^2$. Confirm you obtain (up to scaling) $1,\ x,\ x^2-\tfrac13$, the first three
Legendre polynomials, and that they are mutually orthogonal under this integral.
\end{exercise}

\begin{exercise}[A non-square challenge]
Let $\mat{A}$ be $4\times 3$ with independent columns. Explain why the ``thin'' QR has
$\mat{Q}$ of size $4\times 3$ and $\mat{R}$ of size $3\times 3$, and why
$\mat{Q}\T\mat{Q}=\mat{I}_3$ while $\mat{Q}\mat{Q}\T$ is a $4\times 4$ projector of rank
$3$. What does the rank-$3$ projector do to a vector $\vec{b}\in\R^4$, and how does that
relate to least squares?
\end{exercise}
